\rtmtight
\begin{tblr}{width=\textwidth,colspec={|Q[c,m,2.1cm]|X[l,m]|},hlines,vlines,hspan=minimal,rowsep=1.5pt,colsep=3pt,row{1}={bg=headgray,font=\bfseries}}
\SetCell{c,m}\hfil\logo\hfil & \textbf{POLITEKNIK NEGERI MALANG}\par\textbf{JURUSAN TEKNOLOGI INFORMASI}\par\textbf{PROGRAM STUDI: D4 TEKNIK INFORMATIKA} \\
\SetCell[c=2]{l} \textbf{RENCANA TUGAS MAHASISWA (RTM) DAN RUBRIK PENILAIAN} & \\
Nama Mata Kuliah & Dasar Pemrograman \\
Kode Mata Kuliah & RTI251006 \quad \textbf{sks / Jam perminggu:} 2 SKS/4 jam \quad \textbf{Semester:} 1 \\
Dosen Pengampu & \begin{enumerate}\item Mungki Astiningrum, S.T., M.Kom.\item Imam Fahrur Rozi, S.T., M.T.\item Triana Fatmawati, S.T., M.T.\item Mamluatul Hani'ah, S.Kom., M.Kom.\item Rudy Ariyanto, S.T., M.Cs.\item Noprianto, S.Kom., M.Eng.\end{enumerate} \\
\end{tblr}
\assessmentblock{Tugas Studi Kasus Algoritma (Tugas 1--11)}{Tugas terstruktur dan mandiri}{SCPMK0903-00601, SCPMK0903-00602, SCPMK0704-00603, SCPMK0704-00604}{Mahasiswa menyelesaikan rangkaian tugas studi kasus (Tugas 1--11): analisis permasalahan, version control dan Kanban, tipe data dan operator, pemilihan, perulangan, array, dan fungsi dalam bentuk flowchart/pseudocode.}{Menganalisis permasalahan studi kasus, mengidentifikasi input-proses-output, menyusun algoritma dalam bentuk flowchart/pseudocode, dan mendokumentasikan hasil tugas.}{Lembar jawaban tugas berisi analisis kasus, flowchart/pseudocode, dan dokumentasi penyelesaian.}{Analisis masalah dan konsep dasar (algoritma, version control, Kanban) & 4 & Input-proses-output kasus diidentifikasi lengkap; konsep algoritma, version control, dan Kanban dijelaskan tepat dan diterapkan pada tugas. & Analisis kasus sebagian besar benar, tetapi identifikasi input-proses-output atau penerapan Kanban/version control belum lengkap. & Analisis kasus keliru atau konsep dasar tidak dapat dijelaskan dan tidak diterapkan pada tugas. \\ \\
Ketepatan tipe data, operator, dan sequence dalam algoritma & 2 & Pemilihan tipe data, variabel, operator, dan urutan sequence pada algoritma tepat dan konsisten di seluruh tugas. & Sebagian pemilihan tipe data atau operator kurang tepat, tetapi alur sequence masih benar. & Tipe data, operator, atau sequence banyak yang salah sehingga algoritma tidak berjalan benar. \\ \\
Kebenaran flowchart pemilihan dan perulangan & 8 & Flowchart pemilihan (sederhana/bersarang) dan perulangan (sederhana/bersarang) benar, sesuai notasi, dan menyelesaikan studi kasus. & Flowchart menyelesaikan kasus utama, tetapi terdapat kesalahan notasi atau kondisi pada kasus bersarang. & Flowchart tidak menyelesaikan kasus atau logika pemilihan/perulangan salah. \\ \\
Penerapan array dan fungsi pada studi kasus terintegrasi & 6 & Array 1/2 dimensi dan fungsi (iteratif/rekursif) diterapkan benar pada studi kasus, termasuk kasus terintegrasi, dengan dokumentasi rapi. & Array dan fungsi diterapkan pada kasus dasar, tetapi kasus terintegrasi atau dokumentasinya belum lengkap. & Array atau fungsi tidak diterapkan dengan benar atau tugas terintegrasi tidak dikerjakan. \\}

\assessmentblock{Kuis 1 Konsep Algoritma, Tipe Data, dan Operator}{Kuis (ujian teori closed book)}{SCPMK0903-00601, SCPMK0903-00602}{Mahasiswa mengerjakan kuis pilihan ganda dan/atau essay yang mencakup materi minggu 1--3: konsep dasar algoritma, version control, Kanban, tipe data, variabel, input-output, sequence, dan operator.}{Mengerjakan soal kuis secara mandiri, closed book, dalam waktu yang ditentukan.}{Lembar jawaban kuis (pilihan ganda dan/atau essay).}{Penguasaan konsep algoritma, version control, dan Kanban & 5 & Jawaban menjelaskan pengertian dan pentingnya algoritma, cara kerja version control, dan prinsip Kanban dengan benar sesuai kunci jawaban. & Sebagian besar jawaban benar, tetapi penjelasan konsep masih kurang lengkap atau kurang tepat pada beberapa soal. & Jawaban banyak yang salah atau tidak menunjukkan pemahaman konsep dasar. \\ \\
Ketepatan penerapan tipe data, operator, dan sequence & 5 & Soal penerapan tipe data, variabel, input-output, sequence, dan operator dijawab benar sesuai kunci jawaban. & Sebagian soal penerapan dijawab benar, tetapi terdapat kekeliruan pada operator atau urutan sequence. & Jawaban penerapan banyak yang salah atau tidak sesuai konsep. \\}

\assessmentblock{UTS Algoritma, Sequence, Pemilihan, dan Perulangan}{UTS (ujian teori closed book)}{SCPMK0903-00601, SCPMK0903-00602, SCPMK0704-00603}{Mahasiswa mengerjakan UTS pilihan ganda dan/atau essay yang mencakup materi minggu 1--7: konsep algoritma, tipe data dan operator, sequence, pemilihan sederhana dan bersarang, serta perulangan sederhana.}{Mengerjakan soal ujian secara mandiri, closed book, dalam waktu yang ditentukan.}{Lembar jawaban UTS (pilihan ganda dan/atau essay).}{Ketepatan jawaban konsep dasar algoritma & 6 & Soal konsep algoritma, analisis input-proses-output, version control, dan Kanban dijawab benar sesuai kunci jawaban. & Sebagian soal konsep dijawab benar, tetapi analisis kasus masih kurang tepat. & Jawaban konsep banyak yang salah atau analisis kasus tidak dikerjakan. \\ \\
Ketepatan penyelesaian soal tipe data, operator, dan sequence & 8 & Soal tipe data, variabel, operator, dan sequence dijawab benar termasuk evaluasi ekspresi dan penulisan algoritma. & Sebagian soal dijawab benar, tetapi evaluasi ekspresi atau penulisan algoritma masih keliru. & Jawaban banyak yang salah atau tidak sesuai konsep tipe data dan operator. \\ \\
Kebenaran algoritma pemilihan dan perulangan dalam flowchart & 16 & Soal kasus pemilihan (sederhana/bersarang) dan perulangan sederhana diselesaikan dengan flowchart/algoritma yang benar dan lengkap. & Flowchart/algoritma menyelesaikan sebagian kasus, tetapi terdapat kesalahan kondisi atau alur. & Flowchart/algoritma salah logika atau tidak menyelesaikan kasus ujian. \\}

\assessmentblock{Kuis 2 Array dan Perulangan Bersarang}{Kuis (ujian teori closed book)}{SCPMK0704-00603, SCPMK0704-00604}{Mahasiswa mengerjakan kuis pilihan ganda dan/atau essay yang mencakup materi minggu 9--11: array 1 dimensi, array 2 dimensi, dan perulangan bersarang.}{Mengerjakan soal kuis secara mandiri, closed book, dalam waktu yang ditentukan.}{Lembar jawaban kuis (pilihan ganda dan/atau essay).}{Kebenaran algoritma perulangan bersarang & 5 & Soal kasus perulangan bersarang dijawab dengan algoritma/flowchart yang benar sesuai kunci jawaban. & Algoritma perulangan bersarang sebagian benar, tetapi terdapat kesalahan pada kondisi atau iterasi dalam. & Algoritma perulangan bersarang salah atau tidak dikerjakan. \\ \\
Ketepatan penggunaan array 1 dan 2 dimensi & 5 & Soal deklarasi, akses, dan pengolahan array 1/2 dimensi (searching, sorting, matriks) dijawab benar. & Sebagian soal array dijawab benar, tetapi terdapat kekeliruan indeks atau operasi. & Jawaban array banyak yang salah atau konsep indeks tidak dipahami. \\}

\assessmentblock{UAS Algoritma Terintegrasi}{UAS (ujian teori closed book)}{SCPMK0704-00603, SCPMK0704-00604}{Mahasiswa mengerjakan UAS pilihan ganda dan/atau essay yang mencakup materi minggu 1--16 dengan penekanan pada pemilihan, perulangan, array, dan fungsi dalam penyelesaian studi kasus.}{Mengerjakan soal ujian secara mandiri, closed book, dalam waktu yang ditentukan.}{Lembar jawaban UAS (pilihan ganda dan/atau essay).}{Kebenaran algoritma pemilihan dan perulangan pada soal kasus & 10 & Soal kasus pemilihan dan perulangan (termasuk bersarang) diselesaikan dengan algoritma/flowchart yang benar dan lengkap. & Algoritma menyelesaikan sebagian kasus, tetapi terdapat kesalahan kondisi atau alur. & Algoritma salah logika atau kasus tidak dikerjakan. \\ \\
Ketepatan penyelesaian kasus array 1 dan 2 dimensi & 10 & Soal kasus array (searching, sorting, matriks) diselesaikan dengan deklarasi, indeks, dan operasi yang benar. & Sebagian kasus array diselesaikan benar, tetapi terdapat kekeliruan indeks atau operasi. & Kasus array tidak diselesaikan atau konsep array tidak dipahami. \\ \\
Ketepatan penerapan fungsi iteratif dan rekursif & 10 & Soal kasus fungsi diselesaikan dengan deklarasi, parameter, return value, dan rekursi yang benar. & Fungsi dideklarasikan dan dipanggil benar, tetapi parameter, return value, atau rekursi masih keliru. & Konsep fungsi tidak diterapkan dengan benar atau soal tidak dikerjakan. \\}

\begin{tblr}{width=\textwidth,colspec={|Q[l,m,3.2cm]|X[l,m]|},hlines,vlines,hspan=minimal,row{1-3}={bg=headgray},rowsep=1.5pt,colsep=3pt}
Jadwal Pelaksanaan: & Minggu 1--17 sesuai RPS. Setiap evaluasi memiliki RTM dan rubrik multi-kriteria. \\
Lain-Lain Yang Diperlukan: & LMS, perangkat lunak sesuai mata kuliah (office suite, editor, JDK), dan perangkat presentasi. \\
Pustaka: & Mengacu pustaka utama dan pendukung pada bagian RPS. \\
\end{tblr}
